\section{Conclusion}
Neural Bellman Operators make policy proposal, evaluation, improvement, and certification separate economic computations. That separation is useful when it supports a decision: whether a stored policy remains adequate after a contract changes, whether a portfolio sign is determined by economic incentives rather than numerical error, and whether a reusable continuation approximation pays for its construction.

For the stopped preference economy, the relevant contract contrast is the operating benefit net of the liquidation annuity. Independently evaluated policy features and certified value envelopes provide a uniform welfare guarantee over its full displayed interval. Action-separation bounds identify negative and positive portfolio regions and retain an unresolved switching band. The finite implementation includes every previously diagnosed candidate deviation and repairs the initial neural policies at all dates and states. The exact normalization identity and the separate correlation comparisons identify the operating-duration mechanism without converting it into a universal hedging claim.

The resource experiment establishes why the continuation representation and the actor should be assessed separately. A nonlinear convex continuation fits the tested problem more accurately than a full convex quadratic. Reuse becomes useful at a measured workload, while the actor primarily changes initialization and must earn back its own training cost. The persistent-capacity reference, recursive-utility updates, and temporal-self calculations preserve the broader framework and test its economically distinct operators.

The continuous-time equations, finite transition models, trained proposals, and independent policy values remain distinct objects. This distinction makes the computational results stronger: each numerical guarantee states the policy class, boundary mechanism, parameter range, and error tolerance it actually covers. It also permits subsequent extensions to change the approximation architecture without changing the economic meaning of the certificate.
